\chapter{Échelles, cordes et vingt-six dimensions}

\lectureoverview{Qu’est-ce qui fixe l’espacement du spectre d’une corde ? La réponse est
la tension, une grandeur de dimension masse au carré. Pour les cordes
fondamentales, cette observation conduit aux unités de Planck, à
l’impossibilité pratique de les exciter directement et à l’apparition d’une
forte gravité. Après une pause délibérée, nous revenons à l’état
fondamental de la corde ouverte bosonique. La régularisation de son énergie
de point zéro fournit la voie historique
\emph{facile} vers vingt-quatre directions transverses, et donc vingt-six
dimensions d’espace-temps.}

\section{L’échelle que nous avions laissée de côté}

Jusqu’ici, nous avons parlé de cordes qui vibrent et tournent sans préciser
ce que signifie une unité d’excitation. On peut se représenter un hadron
comme une corde et exciter ses oscillateurs, mais quelle est l’ampleur du
passage d’un état au suivant ? Qu’est-ce qui la détermine ?

Le modèle élémentaire de la corde ne comporte qu’un seul paramètre
ayant une dimension. Pour le trouver, considérons d’abord un ressort
étiré ordinaire. Si ses extrémités sont séparées d’une distance \(\ell\),
la loi de Hooke donne
\begin{equation}
  E_{\mathrm{spring}}=\frac{k\ell^2}{2}.
  \label{eq:l05-hooke}
\end{equation}
Le facteur un demi a été ajouté par la classe après la rédaction de la
première version. Aucun point essentiel de cette estimation initiale ne
dépendra de tels facteurs numériques.

\begin{figure}[H]
  \centering
  \includegraphics[width=0.72\textwidth]{lecture_05_figure_01.png}
  \caption{Le segment étiré et son énergie potentielle élastique selon la loi de Hooke.}
  \lectureframe{00:02:28}
  \label{fig:l05-hooke-board}
\end{figure}

L’énergie dans \eqref{eq:l05-hooke} est non relativiste. Dans la description à
impulsion infinie utilisée dans les cours précédents, l’énergie interne
non relativiste est proportionnelle non pas à \(m\), mais à \(m^2\). Nous
écrivons donc, de manière schématique,
\begin{equation}
  m^2\sim k\ell^2,
  \qquad
  m\sim \sqrt{k}\,\ell.
  \label{eq:l05-mass-length}
\end{equation}
Dans le référentiel au repos de la corde, sa masse est égale à son énergie ;
la seconde relation prend donc la forme
\begin{equation}
  E_{\mathrm{rest}}\sim T_s\ell,
  \qquad
  T_s\equiv\sqrt{k}.
  \label{eq:l05-tension}
\end{equation}
Le coefficient \(T_s\) est la tension de la corde : l’énergie stockée par
unité de longueur.

\begin{figure}[H]
  \centering
  \includegraphics[width=0.78\textwidth]{lecture_05_figure_02.png}
  \caption{Au tableau, on passe de \(m^2\sim k\ell^2\) à une masse proportionnelle
  à la longueur.}
  \lectureframe{00:05:42}
  \label{fig:l05-mass-length-board}
\end{figure}

Cette notation est adaptée à l’analogie avec le ressort utilisée dans le
cours. Dans la notation relativiste usuelle des cordes, on écrit plutôt
\begin{equation}
  T_s=\frac{1}{2\pi\alpha'},
  \label{eq:l05-alpha-prime}
\end{equation}
Ainsi, les accroissements du niveau de la corde produisent des
accroissements de \(m^2\) d’un ordre de grandeur égal à \(1/\alpha'\), donc, de façon équivalente, à la tension.
\editorialnote{Le \(k\) effectif du cours est choisi de sorte que
\(\sqrt{k}\) joue le rôle de la tension. L’énoncé indépendant des conventions
est que \(T_s\) est de dimension masse au carré et que la pente de Regge
vérifie \(\alpha'\propto T_s^{-1}\).}

\subsection{Dimensions et écarts entre niveaux d’oscillation}

Posons \(c=\hbar=1\). L’énergie a alors la dimension de l’inverse d’une longueur :
\begin{equation}
  [E]=L^{-1}.
\end{equation}
La tension étant l’énergie par unité de longueur,
\begin{equation}
  [T_s]=L^{-2}=E^2.
  \label{eq:l05-tension-dimension}
\end{equation}
C’est pourquoi le spectre d’une corde est naturellement espacé en masse au carré.

\begin{figure}[H]
  \centering
  \includegraphics[width=0.68\textwidth]{lecture_05_figure_03.png}
  \caption{Unités naturelles au tableau : \(c=\hbar=1\), l’énergie est l’inverse
  d’une longueur et la fréquence d’oscillation est fixée par l’échelle du ressort.}
  \lectureframe{00:06:05}
  \label{fig:l05-natural-units-board}
\end{figure}

L’oscillateur ordinaire exprime la même idée sous un autre angle. Pour un
ressort de masse \(\mu\),
\begin{equation}
  \omega=\sqrt{\frac{k}{\mu}},
  \qquad
  \Delta E=\hbar\omega.
\end{equation}
Un ressort plus raide a une fréquence plus élevée et un écart quantique plus
grand. Une corde contient une infinité de tels modes normaux, mais leur
échelle commune est fixée par \(T_s\). Dès que la tension est connue,
l’espacement du spectre de la corde l’est aussi.

La tension a aussi la dimension d’une force, puisque le produit d’une
force par une distance est une énergie. Imaginons que l’on fixe une
extrémité et que l’on suspende un poids à l’autre. Une corde relativiste
idéale telle que \(E=T_s\ell\) peut soutenir une force déterminée par \(T_s\),
indépendamment de sa longueur. La tension hadronique déduite des énergies
d’excitation observées correspond, dans cette comparaison délibérément
fantaisiste, au poids d’un camion. On n’a jamais suspendu de camion à un
méson ; le spectre mesuré donne la tension, et la tension permet cette
comparaison mécanique.

Les cordes fondamentales seraient bien plus raides. Si elles sont à
l’origine des électrons, des photons et des gravitons, leurs premières
excitations d’oscillation doivent se situer très au-dessus des énergies des
particules observées. Dans la même expérience gravitationnelle fantaisiste,
leur tension pourrait soutenir un objet de masse galactique. Cette
comparaison ne porte pas sur l’ingénierie. Elle illustre le fait qu’une
tension énorme implique à la fois une énergie d’excitation élevée et une
fluctuation de l’état fondamental très faible.

\section{Des unités construites à partir de constantes universelles}

Quelle échelle est suffisamment élevée ? L’échelle de Planck constitue le
repère naturel. Plutôt que de simplement la citer, construisons-la.

Le mètre, le kilogramme et la seconde sont d’excellentes unités de
laboratoire, mais leur origine est liée à l’échelle humaine. Le mètre est
commode pour mesurer du tissu ou une corde ; le pied porte le nom d’une
partie du corps ; la seconde représente une fraction du jour facile à
appréhender. On pourrait poser que le rayon du proton vaut un et simplifier
certaines formules nucléaires, mais le proton n’est qu’une particule
composite parmi d’autres. Un tel choix ne simplifierait pas toute la physique.

Trois constantes dimensionnées ont une portée plus universelle :
\begin{itemize}
  \item \(c\) fixe la vitesse limite pour tout système matériel ;
  \item \(\hbar\) intervient dans les relations d’incertitude de tout système
  quantique ;
  \item \(G\) couple la gravité à toute forme de masse-énergie.
\end{itemize}
Choisir des unités telles que \(c=\hbar=G=1\) revient à construire des unités
de longueur, de temps et de masse à partir de ces constantes.

Leurs dimensions usuelles sont
\begin{equation}
  \begin{aligned}
    [c]&=LT^{-1},
    & [\hbar]&=ML^2T^{-1},\\
    [G]&=L^3M^{-1}T^{-2}.
  \end{aligned}
  \label{eq:l05-dimensions-constants}
\end{equation}
La relation du milieu découle directement de
\(\Delta x\,\Delta p\gtrsim\hbar\). La dernière découle, par exemple, de
\begin{equation}
  a=\frac{GM}{r^2}.
\end{equation}

\begin{figure}[H]
  \centering
  \begin{minipage}[t]{0.48\textwidth}
    \centering
    \includegraphics[width=\textwidth]{lecture_05_figure_04.png}
    \par\smallskip\footnotesize Unités de \(c\) et \(\hbar\).
  \end{minipage}\hfill
  \begin{minipage}[t]{0.48\textwidth}
    \centering
    \includegraphics[width=\textwidth]{lecture_05_figure_05.png}
    \par\smallskip\footnotesize Unités de \(G\) déduites de l’accélération newtonienne.
  \end{minipage}
  \caption{Les éléments d’analyse dimensionnelle rassemblés en cours.}
  \lectureframe{00:25:43--00:26:40}
  \label{fig:l05-dimensional-analysis-board}
\end{figure}

\subsection{D’abord l’aire de Planck}

Cherchons les exposants \(p,q,r\) tels que
\begin{equation}
  G^p\hbar^q c^r
\end{equation}
ait la dimension \(L^2\). Commencer par l’aire évite d’introduire trop tôt
une racine carrée et annonce aussi le rôle particulier que l’aire jouera
plus tard en gravité. À partir de \eqref{eq:l05-dimensions-constants}, l’élimination de la masse et du
temps donne
\begin{equation}
  -p+q=0,
  \qquad
  -2p-q-r=0.
\end{equation}
L’exposant restant de la longueur doit satisfaire
\begin{equation}
  3p+2q+r=2.
\end{equation}
On obtient donc \(q=p\), \(r=-3p\) et \(p=1\). L’aire de Planck vaut
\begin{equation}
  \ell_P^2=\frac{G\hbar}{c^3},
  \label{eq:l05-planck-area}
\end{equation}
et, par conséquent,
\begin{equation}
  \ell_P=\sqrt{\frac{G\hbar}{c^3}},
  \qquad
  t_P=\frac{\ell_P}{c},
  \qquad
  m_P=\sqrt{\frac{\hbar c}{G}}.
  \label{eq:l05-planck-units}
\end{equation}

\begin{figure}[H]
  \centering
  \includegraphics[width=0.78\textwidth]{lecture_05_figure_06.png}
  \caption{L’argument dimensionnel complet pour l’aire et la longueur de Planck.}
  \lectureframe{00:31:34}
  \label{fig:l05-planck-area-board}
\end{figure}

À un ordre de grandeur près,
\begin{equation}
  \begin{aligned}
    \ell_P&\sim10^{-35}\,\mathrm{m},
    &t_P&\sim10^{-43}\,\mathrm{s},\\
    m_P&\sim10^{-8}\,\mathrm{kg},
    &m_Pc^2&\sim10^{19}\,\mathrm{GeV}.
  \end{aligned}
  \label{eq:l05-planck-numbers}
\end{equation}

\begin{figure}[H]
  \centering
  \includegraphics[width=0.72\textwidth]{lecture_05_figure_07.png}
  \caption{Les trois échelles de Planck inscrites au tableau, à un ordre de grandeur près.}
  \lectureframe{00:35:12}
  \label{fig:l05-planck-values-board}
\end{figure}

La longueur et le temps sont inimaginablement petits, tandis que la
masse est étonnamment macroscopique à l’échelle des particules, à peu près
celle d’un grain de poussière fine. Son énergie de masse au repos est de
l’ordre de quelques gigajoules, comparable à l’énergie chimique d’un
réservoir d’essence. En revanche, l’Univers observable s’étend sur environ
\(10^{60}\) longueurs de Planck et son âge est d’environ
\(10^{60}\) temps de Planck. Ces comparaisons visent à donner une réalité
physique à ces unités, et non à améliorer leur précision numérique.

Si la longueur caractéristique de la corde \(\ell_s=\sqrt{\alpha'}\) est proche de
\(\ell_P\), sa taille quantique est proche de \(\ell_P\), sa période
d’oscillation de \(t_P\) et sa première excitation massive de
\(m_P\). Le cours admet des facteurs de dix, de cent ou davantage ; il
s’agit d’une estimation, pas d’une identité dimensionnelle.

\begin{classroomqa}
\audiencequestion{La masse à l’échelle de Planck correspond-elle à l’augmentation
par rapport à l’état fondamental de la particule ? \lecturetimestamp{00:38:37}}
\lecturerresponse{La grandeur régulièrement espacée est la masse au carré. Si l’état
fondamental est léger, la première excitation de la corde a une valeur de
\(m^2\) supérieure de l’ordre de \(m_P^2\) ; sa masse et l’énergie du
centre de masse nécessaire sont donc de l’ordre de \(m_P\).}
\end{classroomqa}

\section{À quel point les deux échelles doivent-elles être proches ?}

\begin{classroomqa}
\audiencequestion{L’analyse dimensionnelle donne \(\ell_P\), mais pourquoi ce
nombre serait-il la longueur d’une corde ? \lecturetimestamp{00:40:21}}
\lecturerresponse{C’est une hypothèse physique étayée par la gravité et par
l’absence de structure des cordes observée ; elle ne découle pas de la seule
analyse dimensionnelle. La longueur de la corde peut être supérieure à la
longueur de Planck, auquel cas son échelle d’énergie est plus basse. Le
cours soutient seulement qu’elle ne devrait pas être une longueur ordinaire
de la physique des particules et qu’une longueur bien inférieure à celle de
Planck se heurte à l’effondrement gravitationnel.}
\end{classroomqa}

La relation entre les deux échelles dépend en fait du couplage. Dans une
compactification complète, \(\ell_P\), \(\ell_s\), le couplage des cordes et
le volume de compactification interviennent ensemble. Il n’existe pas
d’équation universelle \(\ell_s=\ell_P\). Ce qui subsiste de l’argument du cours,
c’est l’argument d’échelle : les effets des cordes doivent apparaître à très
haute énergie, et on ne peut étudier une localisation sous la longueur de
Planck en négligeant la gravité.
\editorialnote{Les constructions modernes de la théorie des cordes n’identifient
pas nécessairement la longueur de la corde à la longueur de Planck en
quatre dimensions. Leur relation dépend du couplage des cordes et des
données de compactification.}

\begin{classroomqa}
\audiencequestion{Si l’on considère la longueur de Planck comme la plus petite
longueur utile, la masse de Planck est-elle aussi la plus petite masse
possible ? \lecturetimestamp{00:43:27}}
\lecturerresponse{Non. Les particules ordinaires sont bien plus légères. La masse de
Planck est l’échelle à laquelle un objet compact peut avoir un rayon
gravitationnel comparable à sa longueur d’onde quantique. Dans le langage
heuristique du cours, c’est l’échelle du trou noir semi-classique le plus
léger, dont la durée de vie est de l’ordre du temps de Planck ; cette
dernière affirmation ne doit pas être prise pour une description précise
d’un objet pleinement quantique.}
\end{classroomqa}

\subsection{Ce que les couplages en cours d’évolution peuvent nous apprendre}

Des indices indirects montrent que la théorie quantique des champs
conventionnelle reste valable jusqu’à des échelles bien au-delà de la
portée directe des accélérateurs. Les trois constantes de couplage de jauge
du modèle standard sont mesurées aux énergies accessibles, et leur
dépendance en fonction de l’échelle peut être extrapolée. Dans le graphique
schématique du cours, les couplages \(U(1)\), \(SU(2)\) et \(SU(3)\) se
rapprochent à une échelle élevée d’unification, peut-être à quelques ordres
de grandeur sous l’énergie de Planck.

\begin{figure}[H]
  \centering
  \includegraphics[width=0.70\textwidth]{lecture_05_figure_08.png}
  \caption{Graphique schématique des couplages en fonction de l’échelle. Les trois courbes de jauge
  se rapprochent avant que l’intensité gravitationnelle croissante ne devienne
  de l’ordre de un.}
  \lectureframe{00:47:18}
  \label{fig:l05-running-couplings-board}
\end{figure}

Il s’agit d’une extrapolation, pas d’une expérience directe au point
d’intersection. Des particules supplémentaires ou des seuils entre les
énergies de laboratoire et celles d’unification peuvent modifier les
pentes. La conclusion prudente est simplement que la théorie des champs
connue ne révèle aucun signe de sous-structure à l’échelle des cordes aux
distances ordinaires, et qu’une échelle plausible pour les cordes se situe
quelque part dans l’énorme intervalle entre les échelles d’unification
inférée et de Planck.

\begin{classroomqa}
\audiencequestion{Quels sont les trois couplages que l’on extrapole ? \lecturetimestamp{00:48:28}}
\lecturerresponse{Ce sont les couplages de jauge associés à \(U(1)\), \(SU(2)\) et
\(SU(3)\) : les interactions d’hypercharge, faible et forte. Dans cette
classification, l’électromagnétisme résulte du mélange électrofaible ; il
ne constitue pas un quatrième facteur de jauge indépendant.}
\end{classroomqa}

\begin{classroomqa}
\audiencequestion{La constante gravitationnelle \(G\) elle-même augmente-t-elle
lorsqu’on sonde des énergies plus élevées ? \lecturetimestamp{00:49:30}}
\lecturerresponse{\(G\) est dimensionnelle ; prise isolément, ce n’est donc pas la
bonne grandeur à comparer à un couplage de jauge sans dimension. À une
énergie caractéristique \(M\), l’intensité gravitationnelle sans
dimension se comporte schématiquement comme \(GM^2/(\hbar c)\). Cette combinaison
augmente avec l’énergie et devient de l’ordre de un à l’échelle de Planck.}
\end{classroomqa}

Les lois de force rendent la comparaison manifeste :
\begin{equation}
  F_{\mathrm{em}}\sim\frac{e^2}{r^2},
  \qquad
  F_{\mathrm{grav}}\sim\frac{GM^2}{r^2}.
  \label{eq:l05-force-laws}
\end{equation}
Ainsi, en unités naturelles, on compare \(e^2\) non pas à \(G\),
mais à \(GM^2\).

\begin{figure}[H]
  \centering
  \includegraphics[width=0.74\textwidth]{lecture_05_figure_09.png}
  \caption{Lois de force électromagnétique et gravitationnelle utilisées pour définir
  l’intensité gravitationnelle sans dimension.}
  \lectureframe{00:50:50}
  \label{fig:l05-force-laws-board}
\end{figure}

\section{Pourquoi la spectroscopie directe des cordes est hors de portée}

Comment établir expérimentalement qu’un électron est une corde ? Nous
ferions comme en spectroscopie des hadrons : le frapper assez fort pour
exciter ses modes internes et rechercher une série de niveaux dont les
masses suivent le motif propre aux cordes. Si le premier niveau massif se
situe près de \(10^{19}\,\mathrm{GeV}\), il est très au-delà de la portée de tout
accélérateur direct.

L’énergie totale n’est pas le seul obstacle. Il faut la concentrer dans
un paramètre d’impact minuscule. Cela fait intervenir les trous noirs.

\begin{classroomqa}
\audiencequestion{À ces énergies, un accélérateur commencerait-il à produire
de petits trous noirs ? \lecturetimestamp{00:50:57}}
\lecturerresponse{C’est ce que suggère l’argument heuristique lorsque l’énergie du
centre de masse est concentrée à l’intérieur de son rayon gravitationnel.
Un trou noir stellaire se forme lorsque l’énergie potentielle
gravitationnelle comprime une étoile massive. Un trou noir microscopique
se formerait plutôt en concentrant l’énergie cinétique dans une région de
collision minuscule. En gravité à quatre dimensions, l’échelle requise est
celle de Planck.}
\end{classroomqa}

Pour donner un ordre d’idée, le gain en énergie d’un accélérateur linéaire
est à peu près proportionnel à sa longueur. Le cours prend environ deux
miles pour atteindre approximativement \(100\,\mathrm{GeV}\), comme repère simple à
l’échelle du SLAC. Atteindre \(10^{19}\,\mathrm{GeV}\) avec le même gradient multiplierait
la longueur par environ \(10^{17}\) et donnerait une machine de taille
galactique.

Même une telle machine aurait besoin d’une certaine luminosité. Les
faisceaux ne sont pas dirigés avec une précision de longueur de Planck,
pour chaque paire de particules ; on fait se croiser un nombre immense de
particules afin que quelques-unes entrent en collision au paramètre
d’impact voulu. Le produit d’une énergie de Planck par un taux de collision
énorme conduit à une puissance tout aussi absurde. Les estimations sont
volontairement approximatives, mais leur conclusion reste solide : la
spectroscopie directe des cordes à l’échelle de Planck n’est pas une
expérience que l’on puisse raisonnablement concevoir.

Les indices indirects devraient être plus structurels : une théorie qui
produit les particules et les interactions observées à basse énergie avec
assez de précision pour être mise à l’épreuve, ou des effets de la
compactification qui subsistent bien en dessous de l’échelle des cordes.
Le cours ne prétend pas qu’une telle construction décisive soit déjà
établie.

\subsection{Dimensions supplémentaires et limite de localisation}

\begin{classroomqa}
\audiencequestion{Les dimensions supplémentaires devraient-elles avoir elles aussi
des tailles proches de la longueur de Planck ? \lecturetimestamp{00:56:23}}
\lecturerresponse{Dans cette image élémentaire, on s’attend à ce qu’elles soient
extrêmement petites, même si l’argument n’en fixe pas précisément la taille.
Si une dimension compactifiée était assez grande pour être explorée à plus
basse énergie, de nouveaux modes de Kaluza--Klein et des effets
quadridimensionnels auraient déjà modifié la théorie des champs que nous
observons. L’extrapolation qui fonctionne impose donc une borne supérieure
à la taille de ces dimensions.}
\end{classroomqa}

\begin{classroomqa}
\audiencequestion{Pourquoi une corde ou une dimension compactifiée ne pourrait-elle
pas être beaucoup plus petite que \(\ell_P\) ? \lecturetimestamp{00:57:48}}
\lecturerresponse{Pour résoudre une distance plus courte, il faut une impulsion plus
grande. Au-delà de l’échelle de Planck, l’énergie nécessaire à la
localisation a un rayon gravitationnel comparable, puis supérieur, à la
région sondée. L’expérience forme un trou noir avant de pouvoir résoudre
une corde ordinaire plus petite.}
\end{classroomqa}

Le principe d’incertitude fournit l’estimation. Localiser à l’échelle
\(\Delta x\) exige

\begin{classroomqa}
\audiencequestion{Une fois la longueur et le temps de Planck choisis, comment la
masse de Planck est-elle déterminée physiquement ? \lecturetimestamp{00:59:04}}
\lecturerresponse{On peut utiliser soit \(\hbar\), soit \(G\) pour convertir
l’échelle de longueur en énergie. Le principe d’incertitude indique que
localiser deux particules à l’intérieur de \(\ell_P\) exige une impulsion
de l’ordre de \(\hbar/\ell_P\), donc une énergie de l’ordre de \(m_Pc^2\).
Équivalemment, c’est la masse pour laquelle les longueurs quantique et
gravitationnelle deviennent comparables.}
\end{classroomqa}

\begin{equation}
  \Delta p\gtrsim\frac{\hbar}{\Delta x},
  \qquad
  E\sim c\,\Delta p.
  \label{eq:l05-localization}
\end{equation}
Pour \(\Delta x=\ell_P\), cette énergie vaut
\begin{equation}
  E\sim\frac{\hbar c}{\ell_P}=m_Pc^2.
\end{equation}
Au même moment, le rayon de Schwarzschild
\begin{equation}
  r_s\sim\frac{GE}{c^4}
\end{equation}
est de l’ordre de \(\ell_P\), à des constantes numériques près. La
localisation quantique et l’effondrement gravitationnel se rejoignent à la
même échelle.

\begin{figure}[H]
  \centering
  \includegraphics[width=0.74\textwidth]{lecture_05_figure_10.png}
  \caption{Les formules de la longueur et de la masse de Planck reprises dans l’argument sur la localisation.}
  \lectureframe{01:00:25}
  \label{fig:l05-localization-board}
\end{figure}

L’argument qui fixe l’échelle est ainsi achevé.

\section{L’état fondamental doit se situer sous le photon}

Après la pause, deux pistes s’offrent à nous : commencer par les
transformations conformes ou présenter la voie plus ancienne et plus simple
vers la dimension critique. Nous choisissons la seconde, avec une mise en
garde. L’argument simple n’est pas facile et ne donne pas la raison la plus
profonde. Il est néanmoins correct et proche de la découverte historique.

Revenons à la corde ouverte bosonique. L’action d’un opérateur de création
transverse produit deux états de polarisation. En l’absence d’un troisième
état longitudinal, l’excitation doit être un vecteur sans masse. En unités
où un quantum du mode \(1\) augmente \(m^2\) d’une unité,
\begin{equation}
  m_{\mathrm{ground}}^2=-1,
  \qquad
  m_{\mathrm{first}}^2=-1+1=0.
  \label{eq:l05-open-intercept}
\end{equation}
L’état fondamental tachyonique reste un problème de la théorie bosonique ;
nous le reconnaissons et le laissons de côté. La question immédiate est
d’où vient l’unité négative dans \eqref{eq:l05-open-intercept}.

Chaque mode de la corde est un oscillateur harmonique. Pour \(\hbar=1\),
l’énergie du vide vaut
\begin{equation}
  E_{0,n}=\frac{\omega_n}{2}.
\end{equation}
Pour la corde, \(\omega_n=n\), si bien qu’une coordonnée transverse contribue
formellement
\begin{equation}
  E_0^{(1)}
  =\frac12\sum_{n=1}^{\infty}n.
  \label{eq:l05-formal-zero-point}
\end{equation}

\begin{figure}[H]
  \centering
  \includegraphics[width=0.70\textwidth]{lecture_05_figure_11.png}
  \caption{La somme des énergies de point zéro \(\sum n/2\) écrite après le rappel des fréquences des modes.}
  \lectureframe{01:06:18}
  \label{fig:l05-zero-point-board}
\end{figure}

La difficulté apparaît maintenant clairement. Chaque terme de
\eqref{eq:l05-formal-zero-point} est positif et la somme diverge. Comment pourrait-elle fournir
un intercept négatif fini ?

\section{Soustraire le mouvement avant de lire le spectre}

Avant de régulariser la somme, rappelons pourquoi l’énergie interne a été
identifiée à la masse au carré. Pour une particule portant une grande
impulsion longitudinale \(P>0\),
\begin{equation}
  E=\sqrt{P^2+m^2}
   =P+\frac{m^2}{2P}
   +O(P^{-3}).
  \label{eq:l05-large-p-energy}
\end{equation}

\begin{figure}[H]
  \centering
  \includegraphics[width=0.70\textwidth]{lecture_05_figure_12.png}
  \caption{Le développement en grandes valeurs de \(P\) qui sépare l’impulsion globale du terme de masse interne.}
  \lectureframe{01:07:35}
  \label{fig:l05-fast-energy-board}
\end{figure}

Le premier terme décrit le mouvement global sans intérêt le long de l’axe
\(z\). Il est conservé et commun à tous les états internes ayant la même
valeur de \(P\). En le soustrayant puis en redimensionnant, on élimine
la dilatation cinématique du temps :
\begin{equation}
  \begin{aligned}
    E-P&=\frac{m^2}{2P}+O(P^{-3}),\\
    2P(E-P)&=m^2+O(P^{-2}).
  \end{aligned}
  \label{eq:l05-internal-energy}
\end{equation}

L’atome donne une image physique de la situation. Accélérons un atome
d’hydrogène par transformation de Lorentz. Son énergie totale augmente,
mais l’écart entre deux niveaux internes, mesuré selon le temps du
laboratoire, diminue parce que l’horloge de l’atome est dilatée. Le facteur
\(1/P\) dans \eqref{eq:l05-large-p-energy} exprime cette réduction. La multiplication par
\(2P\) rétablit l’écart invariant de masse au carré.

L’énergie brute de l’état fondamental n’a donc pas besoin d’être
littéralement \(-1\). Elle peut comporter une grande contribution
indépendante de l’état, associée à l’énergie longitudinale globale, ainsi
qu’un reste interne fini. Le cours réécrit la relation dans le référentiel
boosté pour faire apparaître les deux contributions, avant de revenir à la
somme des oscillateurs.

\begin{figure}[H]
  \centering
  \includegraphics[width=0.76\textwidth]{lecture_05_figure_13.png}
  \caption{Le terme de masse invariant affiché à côté de la grande contribution longitudinale indépendante de l’état.}
  \lectureframe{01:11:34}
  \label{fig:l05-infinite-piece-board}
\end{figure}

Les énergies du vide infinies sont familières en théorie quantique des
champs. On sépare souvent une constante divergente indépendante de l’état
des différences d’énergie observables. Le cours établit un parallèle avec
l’énergie de point zéro du champ électromagnétique dans la pièce et
qualifie ouvertement la procédure d’insatisfaisante, mais standard. Nous
suivons ici la même approche heuristique : isoler la partie divergente et
vérifier s’il subsiste un reste fini bien défini.

\section{Régularisation de la somme}

Introduisons un régulateur positif \(\epsilon\) :
\begin{equation}
  S_1(\epsilon)
  =\sum_{n=1}^{\infty}n e^{-n\epsilon}.
  \label{eq:l05-regulated-sum}
\end{equation}
Pour tout \(\epsilon>0\), l’exponentielle finit par dominer le facteur \(n\),
de sorte que la somme converge. Faire tendre \(\epsilon\) vers zéro supprime
la coupure.

\begin{figure}[H]
  \centering
  \includegraphics[width=0.78\textwidth]{lecture_05_figure_14.png}
  \caption{Le régulateur exponentiel appliqué terme à terme à \(1+2+3+\cdots\).}
  \lectureframe{01:16:42}
  \label{fig:l05-regulator-board}
\end{figure}

Le facteur \(n\) s’obtient par dérivation :
\begin{equation}
  S_1(\epsilon)
  =
  -\frac{d}{d\epsilon}
  \sum_{n=1}^{\infty}e^{-n\epsilon}.
  \label{eq:l05-derivative-trick}
\end{equation}

\begin{figure}[H]
  \centering
  \includegraphics[width=0.76\textwidth]{lecture_05_figure_15.png}
  \caption{La dérivation par rapport au paramètre de coupure fait apparaître le facteur \(-n\).}
  \lectureframe{01:18:05}
  \label{fig:l05-derivative-board}
\end{figure}

La série restante est géométrique :
\begin{equation}
  S_0(\epsilon)
  \equiv
  \sum_{n=1}^{\infty}e^{-n\epsilon}
  =
  \frac{e^{-\epsilon}}{1-e^{-\epsilon}}
  =
  \frac{1}{e^\epsilon-1}.
  \label{eq:l05-geometric-sum}
\end{equation}

\begin{figure}[H]
  \centering
  \includegraphics[width=0.68\textwidth]{lecture_05_figure_16.png}
  \caption{La série géométrique régularisée sous forme fermée.}
  \lectureframe{01:19:12}
  \label{fig:l05-geometric-board}
\end{figure}

Développons pour \(\epsilon\) petit :
\begin{equation}
  S_0(\epsilon)
  =
  \frac{1}{\epsilon}
  -\frac12
  +\frac{\epsilon}{12}
  -\frac{\epsilon^3}{720}
  +O(\epsilon^5).
  \label{eq:l05-bernoulli-expansion}
\end{equation}
Au tableau, ce résultat est obtenu en développant séparément le numérateur
et le dénominateur, en mettant en facteur une puissance de \(\epsilon\), puis
en développant l’inverse. Ce calcul comporte plusieurs vérifications de
signe faites en cours ; l’expression compacte \eqref{eq:l05-bernoulli-expansion} en est le résultat
mis au propre.

\begin{figure}[H]
  \centering
  \includegraphics[width=0.78\textwidth]{lecture_05_figure_17.png}
  \caption{Développements du numérateur et du dénominateur avant la simplification de l’inverse.}
  \lectureframe{01:23:10}
  \label{fig:l05-expansion-board}
\end{figure}

L’application de \(-d/d\epsilon\) donne
\begin{equation}
  S_1(\epsilon)
  =
  \frac{1}{\epsilon^2}
  -\frac{1}{12}
  +\frac{\epsilon^2}{240}
  +O(\epsilon^4).
  \label{eq:l05-regulated-result}
\end{equation}

\begin{figure}[H]
  \centering
  \includegraphics[width=0.78\textwidth]{lecture_05_figure_18.png}
  \caption{Le terme divergent \(1/\epsilon^2\) et le reste fini \(-1/12\) au terme du calcul avec régulateur.}
  \lectureframe{01:29:48}
  \label{fig:l05-minus-twelve-board}
\end{figure}

On résume souvent ce résultat, sur le ton de la plaisanterie, par
\(1+2+3+\cdots=-1/12\). Il ne s’agit pas d’une convergence ordinaire. L’équation
\eqref{eq:l05-regulated-result} indique que ce régulateur particulier produit un terme divergent
et une partie finie indépendante du régulateur :
\begin{equation}
  \underbrace{\frac{1}{\epsilon^2}}_{\text{divergent}}
  \;+\;
  \underbrace{\left(-\frac{1}{12}\right)}_{\text{finite}}
  \;+\;O(\epsilon^2).
\end{equation}
Le cours absorbe le premier terme dans la contribution indépendante de
l’état déjà présente et conserve le second. Il rappelle aussi la mise en
garde attribuée à Pauli : le seul fait qu’une quantité soit infinie ne la
rend pas nulle. Une théorie cohérente doit expliquer pourquoi cette
soustraction est permise.

En rétablissant le facteur un demi dans \eqref{eq:l05-formal-zero-point}, une coordonnée
transverse contribue
\begin{equation}
  E_0^{(1)}=-\frac{1}{24}.
  \label{eq:l05-one-transverse-zero-point}
\end{equation}

\begin{figure}[H]
  \centering
  \includegraphics[width=0.68\textwidth]{lecture_05_figure_19.png}
  \caption{Le facteur un demi transforme le terme fini \(-1/12\) en \(-1/24\) par coordonnée transverse.}
  \lectureframe{01:33:42}
  \label{fig:l05-minus-twenty-four-board}
\end{figure}

\section{Comptage des directions transverses}

Soit \(D_\perp\) le nombre de coordonnées transverses de la corde. Leur
contribution totale au vide est
\begin{equation}
  E_0^{\mathrm{open}}
  =-\frac{D_\perp}{24}.
  \label{eq:l05-open-vacuum}
\end{equation}
Pour que le premier oscillateur de la corde ouverte fasse passer l’état
fondamental de \(-1\) au niveau sans masse, il faut que
\begin{equation}
  -\frac{D_\perp}{24}=-1.
\end{equation}
Par conséquent,
\begin{equation}
  D_\perp=24.
\end{equation}
Ce sont les directions perpendiculaires à la grande impulsion longitudinale.
En ajoutant cette direction spatiale longitudinale et le temps, on obtient
\begin{equation}
  D=D_\perp+2=26.
  \label{eq:l05-critical-dimension}
\end{equation}

\begin{classroomqa}
\audiencequestion{Quand cet argument est-il entré dans la théorie des cordes ? \lecturetimestamp{01:35:27}}
\lecturerresponse{Il remonte à l’époque des cordes hadroniques, vers 1969 et le
début des années 1970. Les trajectoires de Regge étaient déjà connues
expérimentalement. Lorsque le modèle oscillatoire fut étudié pour lui-même,
on constata que le premier état excité avait deux polarisations vectorielles,
et non trois, ce qui imposait qu’il soit sans masse. L’intercept requis
conduisit alors rapidement à vingt-quatre oscillateurs transverses.}
\end{classroomqa}

À l’époque, cela semblait presque trop étrange pour être crédible. Pourquoi
une somme régularisée du vide devrait-elle connaître la dimension de
l’espace-temps ? Pourtant, tout concordait : le vecteur sans masse,
l’intercept et le nombre d’oscillateurs transverses sélectionnaient tous le
même nombre.

\subsection{La corde fermée confirme le décompte}

Pour atteindre son premier niveau physique sans masse, une corde fermée a
besoin d’un mode se propageant vers la gauche et d’un mode se propageant
vers la droite. Son état fondamental doit donc commencer deux unités en
dessous de zéro.

\begin{classroomqa}
\audiencequestion{Si vingt-quatre directions transverses donnent seulement \(-1\),
une corde fermée doit-elle vivre dans quarante-huit directions transverses ? \lecturetimestamp{01:37:47}}
\lecturerresponse{Non. Les cordes ouvertes et fermées coexistent dans le même espace-temps.
Une corde fermée possède déjà deux fois plus d’oscillateurs : une tour de
modes circule dans chaque direction sur la feuille d’univers. Chaque tour
contribue \(-24/24=-1\) ; ensemble, elles donnent \(-2\).}
\end{classroomqa}

Ainsi,
\begin{equation}
  E_0^{\mathrm{closed}}
  =
  2\left(-\frac{D_\perp}{24}\right)
  =-2
  \qquad (D_\perp=24).
\end{equation}
L’ajout d’un oscillateur de chaque secteur élève l’état de deux unités et
aboutit à une masse au carré nulle, exactement là où le multiplet du graviton
était apparu dans le cours précédent.

Cette concordance rend l’argument plus convaincant, mais elle ne constitue toujours pas l’explication profonde. Dans la quantification moderne, la constante de mise en ordre normal, la fermeture de l’algèbre de Lorentz ou de Virasoro et l’annulation de l’anomalie conforme font partie d’une même structure cohérente. Cette structure fixe à la fois l’intercept et la dimension critique ; le terme résiduel fini dû à la coupure est la voie historique qui y mène.
\editorialnote{Le calcul avec régulateur donne une interprétation heuristique de la constante de mise en ordre normal. La dérivation plus profonde de \(D=26\) exige l’invariance conforme sur la feuille d’univers, ce qui équivaut à l’annulation de l’anomalie conforme et à la cohérence de l’algèbre des contraintes.}

La suite présente cette voie plus profonde. Pour commencer, un peu d’analyse complexe suffit : les coordonnées complexes, les fonctions analytiques et les équations de Cauchy--Riemann. L’image physique annoncée est celle de la pâte à tire turque---tirée dans un sens, étirée dans un autre, remodelée sans détruire la structure essentielle. L’invariance conforme est le fondement mathématique de cette image.